\subsection{Functions, Rates, and the Concept of Dynamic Change}

The study of dynamic systems begins with a fundamental concept from mathematics: the function. A function describes a relationship between two quantities, called the input and the output, such that for each valid input there is exactly one output. When we say that $y$ is a function of $x$, we mean that the value of $y$ is determined by the value of $x$ according to some rule. This rule might be expressed as an algebraic formula such as $y = x^2$, as a geometric relationship such as $y$ being the height of a curve above the point $x$, or as a physical process such as the temperature of a room at time $x$. The key idea is that functions capture how one quantity depends on another.

Consider a concrete example that will accompany us throughout this discussion. Suppose we are tracking the position of a car traveling along a straight highway. If we let $t$ represent time measured in seconds from some starting moment, and let $p(t)$ represent the car's position measured in meters from a fixed reference point, then $p$ is a function of $t$. When $t = 0$, the car might be at position $p(0) = 100$, meaning it begins 100 meters from our reference point. When $t = 5$, the car might have moved to $p(5) = 150$, indicating that in those five seconds it has traveled an additional 50 meters. The function $p(t)$ captures the entire history of the car's motion, telling us exactly where the car is at every instant.

While knowing where the car is at each moment is valuable, dynamic systems are equally interested in how things change. This leads us to the concept of a rate of change. Given any function $f(x)$, the rate at which $f$ changes with respect to $x$ is not constant in general; it varies depending on where we are in the domain. Consider again our car: when traveling on a straight highway with cruise control set to 60 miles per hour, its position changes at a constant rate of 60 miles per hour. But if the driver must slow down for traffic and then accelerate again, the rate of position change varies with time. The mathematical tool that captures this varying rate of change is the derivative.

The derivative of a function $f$ with respect to its input $x$ is itself a new function, denoted $f'(x)$ or $\frac{df}{dx}$, which gives the instantaneous rate of change of $f$ at each point $x$. To build intuition for what the derivative represents, we examine three complementary interpretations: geometric, physical, and algebraic.

\begin{figure}[htbp]
\centering
\begin{tikzpicture}[scale=1.2]
\draw[->, UofSCBlack, thick] (0,0) -- (5,0) node[right] {$x$};
\draw[->, UofSCBlack, thick] (0,0) -- (0,4) node[above] {$f(x)$};

\draw[UofSCGarnet, thick, domain=0.5:4.5, samples=100] plot (\x, {0.5*\x*\x - 1.5*\x + 2});

\draw[UofSCBlack, dashed] (1,0) -- (1,1.5);
\draw[UofSCBlack, dashed] (3,0) -- (3,1.5);

\draw[UofSCBlack] (1,1.5) circle (1.5pt) node[above right] {$A$};
\draw[UofSCBlack] (3,1.5) circle (1.5pt) node[above left] {$B$};

\draw[UofSCBlack, ultra thick] (1,1.5) -- (4.5,3.75);
\draw[UofSCBlack, ultra thick] (0.5,0.375) -- (2,1.875);

\draw[UofSCBlack] (2.5,2.6) node {secant};
\draw[UofSCBlack] (1.4,1.8) node {tangent};

\draw[->, UofSCBlack] (2,0.2) -- (2.8,0.2);
\draw[UofSCBlack] (2.4,0.4) node {$\Delta x$};

\draw[->, UofSCBlack] (3.2,1.5) -- (3.2,2.2);
\draw[UofSCBlack] (3.4,1.85) node {$\Delta y$};

\end{tikzpicture}
\caption{Geometric interpretation of the derivative: the slope of the secant line between two points approaches the slope of the tangent line as the points converge.}
\end{figure}

From a geometric perspective, the derivative at a point represents the slope of the tangent line to the function's graph at that point. To understand this, imagine selecting two distinct points $A$ and $B$ on a curve $y = f(x)$, with coordinates $A = (x_1, f(x_1))$ and $B = (x_2, f(x_2))$. The line passing through these two points is called a secant line, and its slope is given by the familiar rise-over-run formula $\frac{\Delta y}{\Delta x} = \frac{f(x_2) - f(x_1)}{x_2 - x_1}$. This slope represents the average rate of change of $f$ between $x_1$ and $x_2$. Now imagine sliding point $B$ closer and closer to point $A$, so that the distance $\Delta x = x_2 - x_1$ becomes increasingly small. As this distance approaches zero, the secant line rotates and converges to a single line that just touches the curve at point $A$ without crossing through it. This limiting line is the tangent line at $x_1$, and its slope is exactly the derivative $f'(x_1)$.

To make this precise, the derivative is defined as a limit. If we let $h$ represent the small distance between $x_1$ and $x_2$, so that $x_2 = x_1 + h$, then the secant slope becomes $\frac{f(x_1 + h) - f(x_1)}{h}$. The derivative is what this expression approaches as $h$ becomes arbitrarily small:

\[
f'(x_1) = \lim_{h \to 0} \frac{f(x_1 + h) - f(x_1)}{h}.
\]

This limit may not exist for all functions at all points, and those functions for which it does exist everywhere are called differentiable. Throughout our study of dynamic systems, we will work primarily with differentiable functions because their derivatives tell us how the system evolves.

From a physical perspective, the derivative has an immediate and intuitive meaning when our function describes a quantity changing over time. If $p(t)$ represents the position of our car at time $t$, then the derivative $p'(t)$ represents the instantaneous velocity of the car at that moment. Velocity is precisely the rate at which position changes with respect to time. If we measure position in meters and time in seconds, then velocity is measured in meters per second. When the car is moving forward, its velocity is positive; when it reverses direction, its velocity becomes negative. When the car comes to a complete stop at a traffic light, its velocity is zero. The derivative captures all of these physical realities directly.

To see this connection clearly, consider the average velocity of the car between time $t_1$ and time $t_2$. This average is simply the total change in position divided by the total change in time: $v_{\text{avg}} = \frac{p(t_2) - p(t_1)}{t_2 - t_1}$. This is exactly the same formula as the secant slope we encountered in the geometric interpretation. The instantaneous velocity at time $t_1$ is what this average velocity approaches as we make the time interval smaller and smaller, which is precisely the limit definition of the derivative $p'(t_1)$. Thus, the derivative unifies our geometric and physical intuitions: it is both the slope of the tangent line and the instantaneous velocity.

\begin{table}[htbp]
\centering
\begin{tabular}{|c|c|c|c|}
\hline
\rowcolor{UofSCGarnet}
\textcolor{white}{\textbf{Quantity}} & \textcolor{white}{\textbf{Symbol}} & \textcolor{white}{\textbf{Units}} & \textcolor{white}{\textbf{Interpretation}} \\
\hline
Position & $p(t)$ & meters (m) & Where the car is \\
\hline
Velocity & $p'(t)$ & meters per second (m/s) & How fast position changes \\
\hline
Acceleration & $p''(t)$ & meters per second squared (m/s$^2$) & How fast velocity changes \\
\hline
Jerk & $p'''(t)$ & meters per second cubed (m/s$^3$) & How fast acceleration changes \\
\hline
\end{tabular}
\captionsetup{type=table}
\caption{Derivatives of position and their physical interpretations.}
\end{table}

Beyond velocity, we can consider higher-order derivatives. The derivative of velocity with respect to time is acceleration, denoted $a(t) = p''(t)$, which tells us how quickly the car's speed is increasing or decreasing. The derivative of acceleration is called jerk, which measures how abruptly the acceleration is changing. While these higher-order derivatives are important in many applications, the fundamental insight is that each derivative represents a rate of change of the quantity described by the previous derivative.

From an algebraic perspective, the derivative provides a powerful tool for approximation. Near any point where a function is differentiable, the function can be approximated by a linear function whose slope matches the derivative. This is called the linear approximation or tangent line approximation of the function at that point. If we wish to approximate the value of $f(x)$ near a known point $x = a$, where we know both $f(a)$ and $f'(a)$, we use the formula:

\[
f(x) \approx f(a) + f'(a)(x - a).
\]

This formula says that to estimate $f(x)$, we start at the known value $f(a)$ and add the change that would occur if the function continued to change at its instantaneous rate $f'(a)$ over the distance $x - a$. The expression $f'(a)(x - a)$ represents the linear correction based on the derivative.

Let us work through a complete example to see how this approximation works in practice. Consider the function $f(x) = \sqrt{x}$, which gives the square root of $x$. Suppose we know that $\sqrt{4} = 2$, and we wish to estimate $\sqrt{4.1}$, which is close to 4. First, we compute the derivative of $f(x)$. Using the power rule, $f'(x) = \frac{1}{2\sqrt{x}}$. Evaluating at $a = 4$, we find $f'(4) = \frac{1}{2 \cdot 2} = \frac{1}{4} = 0.25$. Now applying the linear approximation formula:

\[
\sqrt{4.1} \approx \sqrt{4} + \frac{1}{4}(4.1 - 4) = 2 + 0.25 \cdot 0.1 = 2 + 0.025 = 2.025.
\]

To verify this approximation, we can compute the exact value. Using a calculator, $\sqrt{4.1} \approx 2.0248457$. Our approximation of $2.025$ is remarkably accurate, with an error of only about $0.00015$. This demonstrates the power of linear approximation: by knowing the function's value and its rate of change at a single point, we can accurately predict nearby values. The closer our target point is to the known point, the better the approximation becomes.

Now we connect these mathematical concepts to the study of dynamic systems. A dynamic system is a mathematical model that describes how a quantity evolves over time. The central insight of dynamic systems is that the rate of change of the system's state depends on the current state itself and possibly on external inputs. In the language of calculus, this means that the derivative of the state variable is a function of the state variable and any inputs.

For our car example, the position $p(t)$ is the state of the system. The rate of change of position, which is velocity $p'(t)$, depends on the current state (the car's speed and direction) and on external inputs (the driver's acceleration or braking). A complete model of the car's motion would express how the velocity changes based on these factors. For instance, if the driver presses the accelerator, velocity increases; if the driver brakes, velocity decreases; if the car encounters a hill, gravity affects both velocity and position.

This relationship between rates of change and current state is the fundamental structure that will govern our entire study of control systems. We will frequently encounter equations of the form:

\[
\frac{dx}{dt} = f(x, u, t),
\]

where $x$ represents the state of the system, $u$ represents external inputs, $t$ represents time, and $f$ is some function that describes how the rate of change depends on these quantities. Solving such equations—determining how the state $x$ evolves given its rate of change—requires the tools of differential equations, which we will develop in subsequent sections.

To solidify this connection, consider a simple dynamic system where a cup of hot coffee cools down over time. Let $T(t)$ represent the temperature of the coffee at time $t$, measured in degrees Celsius above room temperature. Experience tells us that the coffee cools faster when it is hotter and slower as it approaches room temperature. A simple model for this cooling is given by Newton's Law of Cooling:

\[
\frac{dT}{dt} = -kT,
\]

where $k$ is a positive constant that depends on properties of the cup and the coffee. This equation states that the rate of temperature change is proportional to the current temperature above room temperature, with the negative sign indicating that the temperature decreases. Here, the rate of change $\frac{dT}{dt}$ depends only on the current state $T$, with no external input. This is an autonomous dynamic system, and we will learn how to solve such equations to predict the coffee's temperature at any future time.

In summary, functions describe relationships between quantities, derivatives describe how those quantities change, and dynamic systems describe how rates of change depend on current states and inputs. This framework—the language of functions and their derivatives—provides the mathematical foundation for understanding and designing control systems that regulate the behavior of physical, biological, and engineered systems.
